% 本文件由 LaTeX 排版 Agent 根据有效 translation.json 自动生成。
% author=Gregory Thaumaturgus; work_id=0604; title=向奥力振致词并颂扬

\chapter{向\Person{奥力振}致词并颂扬}

% structure: p0001 source=h1 target=omitted reason=page-title-already-rendered-by-parent

% structure: p0002 source=h2 target=section reason=relative-heading-rank; base=h2

\section{论题一：八年来，\Person{格列高利}放弃演说练习，忙于主要研究罗马法与拉丁语}

沉默在许多场合对许多人来说都是一件极好的事，而此刻也尤其适合我，无论有意无意，我都可以闭上嘴唇，感到有必要保持沉默。因为我不善于那些优美典雅的演说——那些以连贯不断的思路、精选的措辞和词语所讲或所写的演说——而且或许我天性不太适合成功培养这门优雅而真正希腊式的艺术。此外，自从我偶然发表或撰写任何长篇或短篇的演讲以来，已经八年了；在此期间，我也没有听到任何人私下发表或撰写任何东西，或公开宣讲任何赞颂性或辩论性的演说，除了那些可敬的人——他们投身于高贵的哲学研究，不太关心语言的优美和表达的风雅。因为，他们只把词语放在次要位置，而是精确地致力于探究和阐明事物本身，按照它们各自的本然。在我看来，他们并不是不希望，而是非常希望，将他们高贵而精确的思想成果用高贵而优美的语言表达出来。然而，也许他们无法如此轻易地运用这种神圣而神性的能力来发挥其自身恰当的概念，同时练习一种措辞雄辩的言说方式，从而在同一个人心中——而且是这渺小的人心中——同时掌握两种成就，这两种成就是两个不同的人的恩赐，并且实际上是彼此最对立的。因为沉默确实是思想和发明的朋友和助手。但如果一个人追求口才和言辞之美，除了通过词语的学习和不断的练习之外，他无法获得它们。此外，另一门学问完全占据了我的头脑，而嘴巴束缚了舌头，如果我想用希腊人的声音发表任何哪怕最简短的演讲；我指的是我们智者们那些可敬的法律，\Place{罗马帝国}所有臣民的事务现在都受这些法律的指导，而这些法律的制定和学习都不容易。这些法律本身是明智而精确的，多种多样而令人钦佩，总而言之，是最彻底希腊式的；它们用\OriginalTerm{拉丁语}{Roman tongue}表达并交给我们，这是一种奇妙而宏伟的语言，非常适合王权，但对我而言却是困难的。即使我说我希望它容易，也无法改变。因为我们的言辞无非是我们心灵状态的一种图像，我们应该允许那些有口才的人，就像一些优秀的艺术家，同样精通艺术并且拥有丰富的色彩一样，拥有自由描绘他们的词语图画，不仅限于单一色调，而且可以有各种描述，并在丰富的花朵混合中拥有最美丽的美，毫无阻碍。\par

% structure: p0004 source=h2 target=section reason=relative-heading-rank; base=h2

\section{论题二：他试在\Person{奥力振}面前谈论其近乎神圣的禀赋，承认自己被引导进入其手中，出乎一切意料}

但我们，如同任何贫乏的人，没有这些多样化的专门知识——或许是因为从未拥有过它们，又或许是因为失去了它们——不得不只使用木炭和瓦片，也就是说，那些粗俗普通的词语和短语；并借助它们，尽我们所能，表达我们心灵的天然性情，用我们现有的语言来表达，努力展示我们心灵形象的印象，即使不清晰或不华丽，至少如木炭画般忠实，欣然欢迎任何可能从各处呈现的优雅而雄辩的表达，尽管我们并不看重它们。此外，还有第三种情况阻碍并劝阻我进行这一尝试，它比其他情况更使我退缩，并劝我无论如何保持沉默——我指的是主题本身，它曾使我对谈论它充满雄心，但现在却使我犹豫和拖延。因为我的目的是谈论一个人，他确实具有人的外貌和名声，但在那些能够沉思他智慧量度之伟大的人看来，他却被赋予了更高贵、近乎神圣的能力。我即将赞美的并非他的出生或身体训练，这使我此刻以过度的谨慎拖延和推迟。同样，也不是他的力量或美貌；因为这些是年轻人的颂词，其表达是否相称并不重要。因为，就转瞬即逝的事物发表演讲，这些事物以多种方式迅速消亡，并以处理大事的宏伟与庄重以及如此怯懦的拖延来谈论它们，似乎是徒劳无益的。诚然，如果有人提议我谈论那些无用且无实质的东西，那些我绝不会自愿考虑谈论的东西——如果我说，有人提议我谈论任何此类性质的东西，我的演讲就不会有这种谨慎或恐惧，以免在任何陈述中显得有失主题的价值。但现在，我的主题涉及那人身上最肖似天主的方面，以及他身上与天主最有密切关系的方面，那确实被限制在这可见的、可朽的形体之内，但却最热切地追求与天主的相似；我的目的是提及这一点，并着手处理更重要的事，并借此向天主性表达我的感恩，因为我有幸遇到这样的人，超乎人的期望——诚然，不仅超乎他人的期望，也超乎我内心的期望，因为我从未在任何时候将此殊荣置于眼前，也未期望过；我说，我的目的，作为一个微不足道、完全缺乏理解的人，要着手处理这样的主题，我退却、犹豫并渴望沉默，并非没有理由。事实上，保持沉默对我来说似乎也是安全的做法，以免在感恩的表达之下，我可能因轻率而用卑贱、琐碎和极其陈腐的言辞谈论高贵而神圣的主题，从而不仅未能达到真理，甚至，在我能力所及，可能对那些相信真理处于此类范畴的人造成伤害，因为一篇软弱无力的论述被撰写出来，更可能引起嘲笑，而非证明其活力与主题的庄严相称。但是，与你相关的一切都超越伤害和嘲笑，哦，亲爱的灵魂；或者，更可以说，神性在此永远保持其本来面目，不为我们琐碎和无价值的言语所动或受损。然而，我不知道我们如何能逃脱鲁莽和轻率的指责，因为我们竟在愚昧中，且几乎缺乏智慧或准备，就试图处理这些重要且可能超出我们能力范围的事。诚然，如果在别处与别人，我们曾渴望在这些事上做出如此年轻般的努力，我们必定会大胆而勇敢；尽管如此，在这种情况下，我们的鲁莽可能不会被归咎于无耻，因为我们并非在你面前大胆尝试。但现在，我们将充分体现愚昧的全部份额，甚至可以说我们已经充分体现，因为我们竟敢以未洗的脚（如俗语所说）进入那些耳朵，圣言本身——并非像一般大众那样裹着脚，且可以说在隐晦难解的言语的厚厚覆盖之下，而是赤着脚（如果可以这样说）——清晰明了地进入其中，并在那里居住；而我们，只能提供这些人类话语中的废物和泥浆，竟敢将之倾入那些习惯于只听神圣纯洁话语的耳朵中。因此，我们已越轨至此，或许已经足够；现在，至少，应该约束自己，不再继续我们的演讲。我确实非常乐意在此停止。然而，既然我已经做了这鲁莽的冒险，如果对我在这事上的冒昧能给予任何宽恕，请允许我首先解释引导我进入这项艰巨事业的原因。\par

% structure: p0006 source=h2 target=section reason=relative-heading-rank; base=h2

\section{论证三．——他被感恩之心的渴望所激励而谈论他。他认为自己应尽最大能力感谢他。一切福佑的起始来自天主；然而无法向他回报相称的感谢。}

在我看来，忘恩负义是一种致命的恶；确实是致命的恶，甚至是最致命的恶。因为当一个人接受了某种恩惠，却未能试图以口头表达谢意作为回报，而其他事情又非他能力所及时，这便表明他要么是个完全无理性的人，要么是缺乏责任意识的人，要么是个毫无记性的人。再者，尽管一个人自然而然地立刻感受到并知晓所领受的恩惠，然而，若他不将这些义务的记忆保持到日后，并向恩主提供某种感恩的证据，这样的人便是迟钝、忘恩负义、不虔诚的家伙；他所犯的过错，无论对伟人还是对小人物都是不可原谅的：—若设想一位伟大而高尚的人不常常满怀感恩与敬意地念及自己所得的伟大恩惠，或设想一个渺小而可鄙的人不尽全力赞美和颂扬他的恩主，不仅是在大事上，也在小事上。因此，伟大的人和那些智力超群的人，由于他们更充裕、更丰富的条件，理当根据他们的能力向恩主献上更大、更相称的赞美。但卑微的人和处境贫乏的人，也不应忽略那些服侍他们的人，不应漫不经心地接受他们的服务，也不应灰心丧气，仿佛自己拿不出任何配得上或完美的东西；相反，他们虽贫穷，却应有良好的情感，不以他们所尊荣者的能力来衡量，只以自己的能力来衡量，他们应当按自己目前力量的尺度向他表示敬意—这样的敬意对受尊敬者而言，很可能令他感激和喜悦，而且在他眼中，其价值不亚于任何伟大而辉煌的奉献，只要奉献是怀着明确的诚意和真诚的心献上的。在圣经中就有这样的记载，\BibleRef{路 21:2}某位贫穷卑微的妇人，当那些富裕有力的人大量而慷慨地捐献自己的财富时，她独自一人投进了一个小钱，甚至是最小的奉献，然而这却是她全部的生活费，并获得了献上最大祭献的见证。因为，依我判断，圣经的话并没有以献出物质的巨大外表为标准来衡量奉献的价值与辉煌，而是以献者的心意和态度为标准。因此，我们绝不可因害怕自己的感恩不足以回报所受的恩惠而逃避这一义务；相反，我们应当勇敢尝试一切，以便献上感恩，即便不是相称的，至少是我们力所能及的，作为应尽的回报。但愿我们的讲论，即使达不到完美的标准，至少能在某种程度上达到目标，避免任何忘恩负义的印象！因为一种持续的沉默，以不能谈论配得上主题的借口而自圆其说，是虚荣且邪恶的；但始终努力作出相称的回报，即使献上感恩之辞的人能力不及主题的功绩，却标志着良好的品格。至于我，即使我不能按主题应得的方式来谈论，我也不会保持沉默；但当我尽了自己一切所能之后，我便可以自慰。那么，就让我这感恩之辞以此方式展开吧。对天主，确实是宇宙的天主，我不敢设想用这样的言辞谈论祂：然而，我们一切福佑的开端都来自祂；因此，我们的感恩、赞美或颂扬的开端也应该是以祂为对象。但事实上，即使我完全投身于这一职责，而且不是像现在这样—即身为渎神且不洁，与一切不圣洁和污秽的邪恶混杂并被玷污—而是真诚、纯洁至极、最真实、最纯朴、不受任何卑污沾染，即使我如此完全投身，并带着新生婴儿般的纯洁去从事这一工作，我自己也无法为万物的主宰和创始者献上任何相称的荣耀和感恩之礼；因为无论是人们单独个别地，还是全体一致地，以同一精神、同一协调的冲动，仿佛所有纯洁之物都汇聚为一，所有与之相异之物也都转向这一服务，都无法以相称于祂的方式来颂扬祂。至论一个人能在多大程度上正确理解祂的化工，并（若可能）相称地谈论祂，那么，就这能力本身而言—这能力并非来自其他任何人，而唯独从祂领受—他不可能找到比拥有这能力本身更值得感恩的事了。\par

% structure: p0008 source=h2 target=section reason=relative-heading-rank; base=h2

\section{四．\Person{圣子}知道如何相称地赞美\Person{圣父}。我们愿在\Person{基督}内、藉着\Person{基督}将感恩献于\Person{圣父}。\Person{额我略}也感谢他的护守天使，因为他由天使引导至\Person{奥利振}。}

但让我们将赞美和颂歌，就是归给万有之君和总监督者、一切福泽的永恒泉源的，交在祂手中；祂在这事上如同在其他一切事上一样，是我们软弱的治疗者，唯独祂能补充所欠缺的；就是交给我们灵魂的护守者和救主，祂的首生\OriginalTerm{圣言}{Word}，万有的创造者和治理者；只有与祂一起，才能为祂自己和众人，无论私下个别地，还是公开集体地，向圣父献上不间断、不止息的感恩。因为祂本身就是真理、智慧和万有之父的能力，并且祂在圣父内，真实而完全地与圣父合而为一，因此不可能因遗忘、不智或任何分离于祂者所有的软弱，而要么无力赞美祂，要么虽有能力却自愿忽略赞美圣父（我们当然不应如此假设）。因为只有祂能最完美地完成对圣父应得的全部赞美；万有之父使祂与自己合而为一，并通过祂几乎使自己的存在客观地圆满，并以完全相称的能力尊崇祂，如同祂自己受尊崇一样；这一地位是万有中现存受造物中第一个也是唯一一个被赋予的，就是这位圣父的\OriginalTerm{独生子}{Only-begotten}，祂在圣父内，并且是\OriginalTerm{天主圣言}{God the Word}；而我们其余人只能表达感恩和虔诚，前提是我们为从圣父而来的一切恩惠，简朴地依赖祂独一者献上供物，从而向万有之源献上合宜的感恩祭；同时也承认，唯一虔诚的途径正是通过祂呈上我们的纪念。因此，为承认那眷顾我们所有人的、不停息的眷顾，无论是大事还是小事，且一直持续至今，愿这篇圣言被接纳为所有感恩和赞美的合宜且永久的表达——我指第一心智本身的、全然完美、活泼且真正活生生的圣言。但让我们的这篇言辞首先被视为一篇感恩辞，献给这位在众人中独一的神圣人物；如果我还想谈论什么超乎此的，特别是那些不可见却更神性的存有，以及那些特别照顾人的存有，那就献给那位通过某种重大决定，从我的童年起就指定我来统治、养育和训练的那位——我指那位从幼年就喂养我的天主圣天使，正如那位天主所爱的圣人所言的，意指他自己特有的天使。虽然他自己，作为卓越者，确实以此术语指定某位足以与他尊严相配的崇高天使（无论是另一位，还是或许就是\OriginalTerm{大谋议的天使}{Angel of the Mighty Counsel}自身，众人的共同救主，因他的成全而成为他自己特有的护守者，我不清楚知道）——我指的是，他承认并赞美某位更高天使为自己的，不论他是谁。但我们除了向众人的共同治理者献上敬意之外，也承认并赞美那位，不论他是谁，曾是我童年时代奇妙的指导者，在过去一切事情上是我仁慈的导师和护守者。因为导师和护守者的这个职务显然既不适合我，也不适合我的任何朋友和亲属；因为我们都是盲目的，看不见眼前的事物，无法判断什么是正确适宜的；而唯独适合那位事先预见一切有利于我们灵魂之事的：\Emph{那位天使，我说}，祂至今仍支持、教导并引导我；除了所有这些其他恩惠之外，还使我与这个人有了联结，这实际上是对我所行的最重要的事。而且，这也是祂为我成就的，尽管我和我所谈论的那个人之间没有种族或血缘的亲属关系，也没有任何其他纽带，也没有任何邻里或乡土的关联；这些却是多数人之间友谊和联合的基础。但简而言之，祂以真正神圣和智慧的先见，将互不相识、陌生、外邦的我们聚集起来，我们彼此被不同的民族、山川和河流彻底分开，如同人与人之间所能隔离的那样，祂就这样成就了这次对我大有裨益的相会；据我判断，从上天的安排中，从我的出生和最早的教育起，就已预先为我准备了这一恩赐。而这如何实现，若要详述，需要很长时间——不仅我如果细致入微地进入整个主题，试图不遗漏任何细节，即便我略过许多事，只打算概述几个最重要之点，也需要很长时间。\par

% structure: p0010 source=h2 target=section reason=relative-heading-rank; base=h2

\section{论点五。——\Person{格列高利}在此叙述其过往生平。说明他生于外教家庭。十四岁时丧父。他致力于修辞学与法律研究。通过眷顾的奇妙引导，他被带到\Person{奥力振}面前。}

我自出生以来最初的教养，是在父母手中；我父家的生活方式是谬误的，是一种我想没有人会期望我们从中获救的方式；我自己作为孩童，没有智慧，又处于迷信的父亲之下，也没有这种希望。随后是父亲的去世和我的孤苦，这或许也是我认识真理的开端。因为那时，我被首次引向救恩和真理的圣言——我无法说明是怎样——是被迫而非自愿。那时我才十四岁，能有何等决断力呢？然而，正是从那时起，这神圣的圣言开始以某种方式眷顾我，恰在众人共有的理性在我内臻于完备之时；是的，那时它首次眷顾了我。虽然我那时对此不甚在意，但如今思之，我认为在我头上可见到神圣而奇妙的天主眷顾的明确标记，这一巧合如此清晰地显明于我的年岁中，它确保了在此年龄之前的一切谬误行为可归咎于年少无知，而圣言不至于白白赐予尚未具备完备理性的灵魂；同时，当灵魂如今具备理性时，虽未赋予神圣而纯洁的理性，却至少不会缺乏与此理性相称的敬畏，而是人的理性与天主的理性在我内同时并开始运作——一个以我无法解释、却自有其能力的方式相助，另一个则接受帮助。当我反思此事时，我同时充满喜乐与恐惧：我确实因天主眷顾的引领而喜乐，却又因敬畏而战栗，唯恐蒙受如此恩惠后，我仍会在任何方面未能达至终点。但我真不知我的话语为何在此事上停留如此之久，我本意是要讲述那将我带到此人面前的奇妙安排（天主眷顾），却又急于以寥寥数语过渡到后续顺序的内容——当然不是以为我能向那如此对待我的那一位献上应得的赞美、感恩或虔敬（因为若以这样的言辞来称呼我们的论述，而所说却无一与主题相称，我们可能显得傲慢），而只是意图呈上可称为平实的叙述或宣认，或任何其他谦卑的称谓。我唯一幸存的家长——我的母亲——认为，在我已受教于那些通常出身教养不差的男孩所受的其他课业之后，还应跟随一位修辞教师，希望我也能成为演说家。于是我便跟随了这样一位教师；那领域中的评判者当时宣称我不久将成为演说家。我自己不知如何判断此事，也不愿如此；因为那时并无此天赋的明显基础，也无任何能引我至此的力量基础。但那神圣的向导和真正的守护者，总是如此警醒，当我的亲友们不曾考虑此事，我自己也不渴望时，他来了，向我的其中一位老师——我受命跟随他学习罗马语，并非指望我能完全掌握该语言，只求不致全然无知——提出了（扩展学业）的建议；这位老师碰巧也并非全然不谙法律。因此，他将这念头放入这位老师心中，使我通过他深入学习罗马法律。那人热心地承担起对我的教导；我方面，也投身其中——更多是为了取悦那人，而非我自己喜爱此学。当他收我为徒后，便开始满怀热忱地教我。他说了一句话，在我看来是他所有话语中最真实的一句，即我在法律方面的教育将是我最大的\OriginalTerm{盘缠}{viaticum}——他是这样措辞的——无论我是渴望成为法庭上争辩的演说家，还是愿属于另一阶层。他这样表述，本意仅指人间之事；但在我看来，他是受比他自己所想的更神圣的冲动驱使而说出此言。因为，当我无论情愿与否，已精熟这些法律之时，仿佛绳索立时束缚了我的行动，而我从贝鲁特城启程前往此地的原因和机缘出现了：贝鲁特城离此区域不远，有些拉丁化，并被视为这些法律研究的学校。而这位可敬的人来自埃及亚历山大城，他先前恰好居住在那里，因其他情况而迁居此地，仿佛专为与我们会面。至于我，无法解释这些事件的原因，也乐意略过。但这一点是确定的：那时我学习我们法律的意图尚未提供我前来此地与此人会面的必要机缘，因为我也能前往罗马城本身。那么，此事如何成就呢？当时的巴勒斯坦总督突然夺走了我的一位亲友，即我姐妹的丈夫，使他与妻子分离，并违背其意愿将他带至此地，以获得他的帮助，并让他参与治理国家的劳作；因他精通法律，或许至今仍如此\Emph{任用}。然而，他与他同行后，不久便有幸得以接来妻子，重新得到她——他曾违背意愿、满怀苦恼地被与她分离。如此，他碰巧也带着我们与她一同来到同一地方。因我们当时打算旅行，不知去何处，但肯定去任何其他地方而非此地，突然一名士兵出现，携带着指令信函，要我们护送并保护我们的姐妹，使其回到丈夫身边，并让我们自己作她的旅途伴侣；这样我们有机会帮助我们的亲属，尤其是我们的姐妹（使她不必以不得体或充满恐惧犹豫的方式踏上旅途），同时我们的其他亲友也赞同此事，认为这并非毫无益处，因我们可以借此访问贝鲁特城，并在那里全心投入法律研究。因此，一切事情都促使我前往那里——对姐妹的责任感、我自己的学业，此外还有那士兵（也应提及此事），他拥有比案件所需更多的公用车辆，以及比仅需供给我姐妹更多的支票。这些是我们旅途的显见理由；但隐秘且更真实的理由则是：与这人共处的机缘、通过那人接受关于圣言的真理的教导，以及我灵魂为得救恩的益处。这些是将我们带至此地的真正原因，我们对于获得救恩的方式盲目无知。因此，并非那士兵，而是一位神圣的同伴、仁慈的引领者和守护者，始终引领我们平安度过今世整个生命，如同一次长途旅行，带领我们经过其他地方，特别是贝鲁特——那时我们似乎最急于抵达该城——并将我们带至此地安顿，安排并指引一切，直到他终将我们与此人联结，此人将成为我们大部分恩惠的根源。那以如此方式前来的神圣天使，将此职责交托给他，并（如果可以这样说）或许在此安歇——并非因劳苦或任何疲惫（因那些神圣的侍从的类体不知疲倦），而是因他已将我们交托给一位能尽心履行其能力范围内全部照顾和守护之责的人。\par

% structure: p0012 source=h2 target=section reason=relative-heading-rank; base=h2

\section{论点六——\Person{奥力振}用以挽留\Person{额我略}及其弟\Person{阿忒诺多鲁斯}几乎违其所愿的方法；以及俘获二人的爱情。论哲学为虔敬之基，以此旨趣，\Person{额我略}愿全然投身哲学研究，舍弃祖国、父母、法学研习及其他一切学问。论灵魂为自由之本。高贵者不愿与低贱者结合，低贱者却愿与高贵者结合。}

从他接纳我们的第一天起（那日对我来说实在是第一天，是万日中最珍贵的一天，如果可以这样说，因为那时真正的太阳才第一次开始在我身上升起），当我们像田野里的野兽，或像鱼，或像某种落入圈套或罗网拼命想要挣脱逃走的鸟一样想要离开他，逃往贝鲁特或我们家乡时，他却想方设法要使我们紧紧跟随他，运用各种论证，使出所有的解数（正如谚语所说），调动全部力量来达到这一目的。为此，他以大量的赞美和许多高尚的话语称赞热爱哲学的人，宣称只有那些立志过正直的生活、首先认识自己——他们是什么样的人——然后认识真正善的事物——人应当追求什么——以及真正恶的事物——人应当逃避什么——的人，才活得与理性受造物相称。接着他斥责无知和一切无知的人：这样的人很多，像笨重的牲畜那样心智盲目，甚至不理解自己是什么，迷失了方向就像完全缺乏理性一样，既不知道什么是善什么是恶，也不在乎从别人那里学习，反而狂热地追求财富、荣耀和世俗的荣誉以及身体的舒适，为这些事如痴如狂，好像它们才是真正的善。仿佛这些对象很有价值，甚至比一切都宝贵，他们珍视这些事物本身，以及获取它们的技艺，还有使它们得以实现的不同生活道路——即军人、律师和法律研究。他用诚恳而睿智的话语告诉我们，这些就是使我们软弱的东西，因为我们轻视那本当在我们内作真正主宰的理性。我此刻无法一一叙述他为了说服我们踏上哲学之路而对我们发表的那些讲话。他不仅第一天这样对待我们，而且许多天，事实上我们起初习惯于去找他的每一天都是如此；我们从第一次听他讲话时起就被他的论证像箭一样刺穿（因为他罕见地结合了甜美的优雅与说服力以及一种奇异的强制力），尽管我们仍然犹豫不决，在心中与己争论，虽然坚持追求哲学，却没有完全被说服，而不知怎的又完全无法从中抽身，因此总是被他推理的力量像某种更高的必然性一样吸引过去。因为他进一步断言，凡轻视哲学这份恩赐的人，就不可能有对万有之主的真正虔诚——这份恩赐是世上所有受造物中唯独人类被认为配得上并值得拥有的，而所有尚未因疯狂的头脑错乱而完全丧失思考能力的人，不论智慧与否，都会合理地接受它。因此，正如我所说的，他断言（严格地说）不进行哲学思考的人不可能真正虔诚。他这样继续对待我们，直到通过向我们灌输一个又一个这样的论证，最后凭着某种神圣的力量，像用他的推理俘获俘虏一般把我们完全夺去，并把我们安置在（哲学实践）中，使我们在他的技艺面前像无法动弹一样。此外，友谊的刺激也施加在我们身上——这种刺激实在难以抗拒，敏锐而极为有效——这就是他那亲切友善的性情在他与我们谈话和交往时所表现出来的仁慈话语。因为他不仅仅是想用某种推理来笼络我们；他的愿望是以仁慈、真挚和极其善良的心拯救我们，使我们分享哲学所带来的福乐，尤其分享天主在哲学之外赐予他的那些恩赐——这些恩赐超过大多数人，或者可以说超过我们这个时代的所有人。我指的是教导我们虔诚的能力、救恩的圣言，它临到许多人身上并征服所有它临到的人：因为没有任何事物能抵抗它，因为它现在是、将来也永远是万有之王；尽管它仍隐藏着，未被普通人轻易或艰难地认出，以致当他们被问及时，不能明智地谈论它。于是，就像一点火花落在我们内心深处，爱被点燃并爆发成火焰——这爱既是对至圣圣言的爱，祂以不可言喻的美吸引一切人不可抗拒地归向祂，也是对祂的朋友和辩护者——这个人——的爱。我被这爱强烈地击中了，于是被说服放弃了所有那些我们以为合适的事物或追求，其中包括我自豪的法学，——甚至我的祖国和朋友们，无论是当时与我在一起的，还是我已离开的。在我心中只存有一个可爱且值得向往的对象——那就是哲学，以及那位哲学大师，这位受神感的人。\Quote{约纳堂的灵魂与达味的灵魂紧密相连。} \BibleRef{撒上 18:1} 这句话，我后来才在圣经中读到；但在此之前，我就已感受到它，不亚于经上书写的那样：因为事实上，它通过最清晰的启示在那时临到了我。因为不仅仅是约纳堂与达味紧密相连；而是人与人里面那些起统治作用的部分——他们的灵魂——紧密相连；即使人身上所有明显和表露的东西都被切断，当它们自己不愿意时，也无法用任何技巧强迫分离。因为灵魂是自由的，无法被任何方式强迫，即使把它关在秘密牢狱里看守。因为无论理智在哪里，灵魂也按照自己的本性和第一理性在那里。如果它在你看来好像在一个牢狱里，那是通过第二理性向你呈现的。尽管如此，它决不能被阻止按照自己的决定在任何地方存在；反而，它只能在且只有在那些唯独属于它的活动进行的地方和关联中存在，并且被合理地认为是如此。因此，我所经历的在这短短的一句话中表述得十分清楚：\Quote{约纳堂的灵魂与达味的灵魂紧密相连；}这些对象，如我所说，无法被任何方式强迫分离，它们出于自愿当然不会轻易选择分离。我认为，解除这种神圣之爱纽带的力量不在于那较卑微者——他易变、主意多变，单凭自己根本没有结合的容量——而在于那更尊贵者——他坚定、不易动摇，通过他才能系上这些纽带、系上这种神圣之结。因此，不是达味的灵魂被圣言与约纳堂的灵魂联系在一起；相反，后者——即较卑微者的灵魂——受到如此影响而与达味的灵魂相连。因为更尊贵者不会选择与较卑微者结合，因为它自足；但较卑微者，因需要更尊贵者所能给予的帮助，理应与更尊贵者相连，并依附于它：这样，后者仍然保持自身的自足，不会因与较卑微者的连接而受损；而本来无序的较卑微者，在与更尊贵者结合并和谐地连接后，就能毫无损害地通过这种纽带的约束完全服从于更尊贵者。因此，施加纽带是属于更尊贵者的事，而非较卑微者的事；但与对方相连则是较卑微者的事，并且是以这样一种方式，即它没有能力使自己从这些纽带中解脱。于是，凭着类似的约束，我们这位达味曾将我们束于自身；他现在仍然持有我们，并且从那时起一直持有我们，以至于即使我们想要，也无法从他那里解脱。因此，即使我们离去，他也不会释放我的这个灵魂——圣经说，这个灵魂被他如此紧密地与自己相连。\par

% structure: p0014 source=h2 target=section reason=relative-heading-rank; base=h2

\section{论证七——\Person{奥力振}为\Person{额我略}与\Person{阿特诺多鲁斯}预备哲学的精妙技巧。二人之智首先在逻辑学中操练，而单纯关注文字被鄙弃。}

但他在一开始就这样俘虏了我们，并仿佛从四面把我们围困起来，当他把最好的事成就了，当我们觉得与他同住一段时间是好的，那时他就接手管理我们，如同一个熟练的农夫接手一块未曾耕作的田地，完全贫瘠、酸涩、烧焦、坚硬如石、粗糙，或者，也许不是完全荒芜或不结果子，反而可能是天然非常肥沃，只是那时荒废、被忽略、僵硬、因荆棘和野灌木而难以处理；或者如同一个园丁接手一株野生的植物，不产生栽培的果实，虽然可能并非毫无价值，发现它如此后，凭借他的园艺技巧，带来一个栽培的嫩枝并嫁接上去，在中间切开一道裂口，然后将两者接合，捆绑在一起，直到两者的汁液合流，以同样的养分一同生长：因为人们常常看到一棵混杂无价值的树，尽管过去不结果子，却能变得多产，并在野生的根上结出好橄榄树的果实；或者看到一株野生植物因细心园丁的技巧而免于完全无益；或者再看到一株本是栽培且多产的植物，但因缺乏熟练的照料，被留下不加修剪、不浇水、荒废，因此被随意生长的多余枝条所窒息，却仍被引导去完成其发芽的正当功能，并结出先前因多余生长而受阻的果实。就是这样，以这样的态度，他起初接待了我们；他以农夫的技巧审视我们，彻底衡量我们，不把他的注意局限于众目所睹、公开显露的那些事，而是更深入地渗透我们，探查我们最内在的部分，他询问我们，向我们提出建议，倾听我们的回答；每当他在我们身上发现并非完全无果、无益或荒废的东西，他就着手清理土壤，翻耕、灌溉，使一切运转，将他的全部技巧和关心施加于我们，作用于我们的心智。至于荆棘和蒺藜，以及我们的心智（如此不规则且轻率地行动）在其未开化的繁茂和天然的野性中所产生的一切野生草本或植物，他通过驳斥和禁止的过程将其割除并彻底清除；有时以真正的\Person{苏格拉底}方式攻击我们，有时又通过他的论证使我们跌倒，每当他看到我们在他面前变得倔强，像许多未驯服的马匹，跳出赛道，疯狂地乱跑，直到他以一种奇异的说服力和约束力，用他的话语使我们在他之下安静下来，这话语像马嚼子一样在我们口中。起初这对我们来说是一个不愉快的处境，且不无痛苦，因为他对待那些不习惯、仍完全未经训练去服从理性的人，当他用他的论证不断追问我们时；然而他通过这些论证净化了我们。当他使我们变得适应，并成功为我们接受真理的话语做好准备后，更进一步，好像我们现在是一块耕透、松软、准备好使撒在里面的种子生长的土壤，他慷慨地对待我们，按时播种好种子，并注意良好耕作的所有其他关心，各按适当的时候。每当他察觉到我们心智中有任何软弱或卑劣的成分（无论是天性如此，还是因身体过度滋养而变得如此粗俗），他就用他的话语刺伤它们，并用那些精巧的言辞和推理的转折来削减它们，这些言辞和转折起初虽然最简单，但逐渐一个接一个地演变，巧妙地展开，直到达到一种几乎无法掌握或解释的复杂性，并使我们仿佛从睡眠中惊醒，教导我们总是紧握眼前之事的艺术，既不因长度也不因微妙性而失足。如果我们身上有任何轻率冒失的倾向，无论是对于所遇到的一切都赞同，无论对象性质如何，甚至证明是假的，或者是经常反对其他事物，即使它们是真实地说的——他也通过已经提到的那些精巧推理以及其他类似推理（因为哲学的这一分支有多种形式）使我们在这一方面受纪律，并使我们习惯于不只随意、凭偶然冲动时而给予证词，时而又拒绝，而是在仔细审查之后才给予，不仅审查明显的事，也审查隐秘的事。因为许多本身声誉很高、外表可敬的事物，通过悦耳的话语进入我们耳中，仿佛它们是真实的，然而它们却是空洞虚假的，并夺走并占有了我们手中真理的赞同票，然后，没过多久，就被发现是败坏、不值得信任、欺骗性地穿着真理外衣的；并因此太轻易地暴露我们为可笑地受骗的人，轻率地作证给那些绝不应得到赞同票的事物。相反，其他事物是真正可敬、绝非欺骗的，但没有用貌似合理的陈述表达，因此显得似是而非、最不可信，并且因其自身表现而被当作虚假拒绝、被不应当地嘲笑，后来经过仔细调查和检验，被发现是万物中最真实的、完全无可争辩的，尽管一时被鄙视并被认为是假的。因此，他不只是处理明显突出的事物（这些有时是欺骗性和诡辩性的），而且也教导我们探查我们内在的事物，并逐一检验它们，免得其中任何一个发出空洞的声音，并指导我们首先确保这些内在事物，然后才训练我们如此这般地对外在事物表示同意，并对所有这些分别表达我们的意见。这样，我们心智中批判性地处理言辞和推理的能力，就以理性的方式被教育；不是依照著名修辞学家的判断——无论希腊或外国的荣誉与那个头衔相关——因为他们的学问价值不大且非必要：而是依照对所有人最需要的东西，无论是希腊人或外国人，无论智识或文盲，最后，为了避免分别列举每个职业和追求而说得太长，是依照对所有人类最不可或缺的东西，无论他们选择了何种生活方式，如果所有必须彼此谈论任何话题的人都真正关心并感兴趣于免受欺骗。\par

% structure: p0016 source=h2 target=section reason=relative-heading-rank; base=h2

\section{论述八：——然后他依次教导他们物理学、几何学和天文学}

他并不只是把精力局限在那种由辩证法来规范的思维形式里；他也着手处理心智的那种谦卑的官能，即当我们对世界的宏大、奇妙、壮丽且绝对智慧的构造感到惊叹时，以无理性的方式惊讶、被恐惧压倒、像无理生物一样不知如何下结论的官能。他也通过自然科学的其他研究来唤醒和纠正这种官能，阐明并区分受造物的各个不同部分，以令人赞叹的清晰度将它们还原到最初的元素，在论述中清晰地将其全部呈现，遍历整个自然以及每一部分的自然，讨论世间事物的千变万化和变迁，直到他以清晰的教导完全带领我们同行；通过那些他部分从他人所学、部分自己发现的推理，他以对宇宙神圣安排和万物无可指摘的构造的理性惊叹，取代了我们无理性的惊叹。这就是自然哲学所教导的那种崇高而属天的学问——一门最吸引所有人的科学。现在何须再谈论神圣的数学，即几何学——对所有人珍贵且无可争议，以及天文学——其轨迹在高天？这些不同的学问他印在我们的理解中，训练我们，或唤起我们的记忆，或与我们做其它我不知道如何正确称呼的事。他将几何学作为一切不变的根基和稳固的基础清晰地呈现；通过天文学，他将我们提升到最高处之上，并借助这两门科学如同登天的阶梯，使天空对我们成为可行的。\par

% structure: p0018 source=h2 target=section reason=relative-heading-rank; base=h2

\section{论述九：——但最重要的是，他以伦理科学充满他们的心灵；并且他不限于言谈讨论德行，而是以行动证实他的教导。}

此外，关于那些最为重要的事，以及为了它们，整个哲学劳作的家庭（好比采集其他一切学问和长久实践哲学所结出的各种善果）——我指的是关乎道德本性的神圣德行，心灵的各种冲动借以获得均衡而稳定的存续——他也通过这些德行，致力于使我们真正抵御一切患难所带来的忧伤和不安，并赋予我们一种有良好修养、坚定而虔诚的精神，好使我们在一切事上确实有福。他为此辛勤劳作，通过相关的话语（充满智慧和抚慰的倾向），并且常常通过最有力的讲论来触动我们的道德倾向和生活方式。他不仅通过言语，也通过行为，在一定程度上调节我们的倾向——也就是说，通过那种对心灵冲动和情感的静观与观察，凭借其结果，心灵尤其常从纷乱状态归于正直，从混乱状态恢复为判断与秩序的状态。如此，当心灵如同在镜子中观看自己时（我具体地说，一方面看到它内在的罪恶的起源和根源，以及一切不合理的东西，一切荒谬的情感和激情由此而生；另一方面，看到真正卓越和合理的一切，在这些东西的支配下，它自身保持免受伤害和扰乱；然后仔细审视在它内所发现的这些东西），它就能驱逐那些来自低劣部分的东西，那些因纵欲而耗损我们力量，或因沮丧而阻碍和窒息我们力量的东西——诸如快乐和情欲，或痛苦和恐惧，以及伴随这些不同邪恶种类的全部祸患。我说，它能够驱逐并除掉它们，在它们刚刚开始生长时就对付它们，不让它们因片刻延缓而增强，而是立即摧毁并根除它们；同时，它也培育那些真正良善的、来自高贵部分的东西，通过在其初生时养育它们，并小心守护直到它们成熟。因为（他常说）天上的德行就是这样在灵魂中成熟的：即智德，它首先能够从事情本身以及从对我们身外之事（无论善恶）的知识中，立即判断出心灵中的那些动态；节德，它能在这些动态的起初做出正确选择的力量；义德，它将公平分配给每件事；以及那保存一切德行的——勇德。因此，他并不使我们习惯于仅仅在言语上宣称，例如智德是对善恶或当做不当做的知识：因为如果只有教义而无行动，那确实是虚妄无益的学问；那种智德也毫无价值，它不做当做的事，也不使人远离不当做的事，却只能为拥有它的人提供这些知识——尽管我们观察到许多这样的人。同样，他也不满足于仅仅断言节德就是关于应当选择什么和不应当选择什么的知识；虽然其他哲学学派甚至连这一点也不教导，尤其是近代的学派，他们在言辞上如此有力而充满活力（以至于我常常对他们感到惊讶，当他们试图证明天主和人具有同一种德行，而且特别在世上，智者与天主同等），却无法传授智德的真理，使人按照智德的指示行事，也无法传授节德的真理，使人选择他所学到的；同样，他们对于义德和勇德的论述也是如此。然而，这位老师并非仅仅以言语向我们传授有关德行的真理；而是更多地激励我们实践德行，以他所行的榜样更甚于他所教的道理来激励我们。\par

% structure: p0020 source=h2 target=section reason=relative-heading-rank; base=h2

\section{论证十：由此驳斥只说不做的空言智者。}

现在，我恳求当代的哲学家们，无论是那些我本人亲自认识的人，还是那些我从他人传闻中所知的人，我也恳求所有其他的人，善意地接受我刚才所作的陈述。但愿没有人以为我这样说是出于对那人单纯的友谊，或出于对其他哲学家的憎恨；因为如果有人倾向于因他们的言论而钦佩他们，并希望称赞他们，且乐于听到别人对他们作出最光荣的提及，那么我自己就是这样的人。然而，我提及的那些事实具有这样的性质，以至于给哲学本身的名声带来了几乎来自大众的极端嘲笑；我几乎可以说，我宁愿对哲学全然无知，也不愿学习这些人所宣称的任何东西，而我认为不再与这些人在这今生交往是好的——尽管在那一点上，或许我判断有误。但我说，没有人应该认为我作出这些陈述仅仅出于对赞美此人的热忱，或出于对任何其他哲学家的敌意。但请所有人确信，我所说的甚至比他的功绩应得的还要少，以免我显得沉溺于奉承；而且我并不追求雕琢的词语和短语，以及巧妙的赞美手段——我即使在年轻时，在学习公众演讲术教授所教的流行演讲风格时，也从未凭自己的意愿说出过一句不真诚的赞美或对任何人作出过不真实的颂词。因此，在现在这个场合，我也认为不对，若我给自己设定单纯称赞他的任务，却以贬低他人为代价来抬高他。而且，说实话，如果我为了有更伟大的东西可以谈论他，而将他有福的生活与他人的缺点相比较，那只会对他造成伤害。然而，我们并非如此愚蠢。但我要简单地见证我亲身经历的事，摒弃所有不当的比较和言辞上的诡计。\par

% structure: p0022 source=h2 target=section reason=relative-heading-rank; base=h2

\section{论点十一——奥力振是第一位也是唯一一位劝勉额我略在其学识之外加上哲学研究，并以某种方式在他自己身上提供榜样者。论正义、明智、节德与刚毅。格言：认识你自己。}

他也是第一个也是唯一一个促使我学习希腊哲学的人，他以自身的道德榜样说服我听从并持守道德学说，然而那时我尚未被其他哲学家说服（我再次承认这一点）——确实，不是以正当的方式，而是不幸的，对我来说，毫不夸张地说。然而，起初我并没有与许多人交往，只与少数自称教师的人结交，但事实上，他们都仅仅停留在言辞上建立哲学。然而，这个人（他）是第一个通过他的言辞诱导我进行哲学思考的人，因为他在言语教导之前先以行为作为劝诫，并且不仅仅是背诵一些精心研究的句子；不仅如此，他甚至认为根本不应该谈论这个话题，而是以真诚的心志，热切地努力实践所说的事情，并始终努力在品格上显示出他在讲论中所描述的那位将过高尚生活的人，我可以说，他总是展现出智者的典范。但由于我们的论述一开始就打算探讨真理，而非虚浮的言辞，我现在不会称他为智者的典范。然而，如果我选择这样描述他，我也不会远离真理。尽管如此，我现在暂且不提。我不会说他是一个完美的典范，而是说他热切地渴望模仿完美的典范，并以热情和诚挚追求它，甚至超出人的能力，如果我可以这样表达的话；他还努力使我们这些如此不同的人变得像他一样，不仅是关于灵魂冲动的干瘪学说的掌握者和领悟者，而且是这些冲动本身的掌握者和领悟者。因为他推动我们走向行动和学说，并通过同样的观点和方法带领我们，不仅是进入每种德性的一个小部分，而是进入整体，如果我们能够接受的话。他还强迫我们，如果我可以这样说的话，基于灵魂本身的个人行动践行正义，他说服我们研究这一点，将我们从生活的多事焦虑和法庭的喧嚣中拉开，提升我们到更高贵的使命：反省自己，并真实地处理关乎我们自身的事。现在，这就是践行正义，这就是真正的正义，一些古代哲学家也断言过（我认为他们称之为\Emph{个人行动}），并肯定这对于有福来说更有利，无论是对人自己还是对他们身边的人，如果这种德性的确在于按功过回报以及各得其所。因为还有什么其他事物能如此特别属于灵魂？或者什么能如此值得灵魂，如同关怀自身，不向外观看，不忙碌于无关之事，或者简言之，不为自己做最不义的亏待，而是将注意力转向内在，给予自己应得的，并由此正义地行动？因此，他以这种方式勉励我们践行正义，如果我可以这样说的话，几乎是强制性地。他同样教育我们明智——教导我们与自己相处，渴望并努力认识自己，这的确是哲学的最高成就，也被视为最富有预言性的灵作为智慧的最高论证——那条诫命：\Quote{认识你自己}。这正是明智的真正功能，也是属天的明智，古人很好地断言了这一点；因为其中有一种人与天主共通的德性；而灵魂在观看自身如同在镜子中得到锻炼，并在自身中映照出神圣的心智，如果它配得上这样的关系，并为达成某种\OriginalTerm{神化}{apotheosis}而探索出一种不可言说的方法。与此相应的是节制和刚毅的德性：节制，确实保守这种明智本身，这种明智必须在认识自身的灵魂中，如果这是它的命运（因为这种节制再次肯定只是健全的明智）；刚毅，在于坚定地持守所有已说过的职责，不因自愿或任何强制而偏离它们，并持守和坚持所有已奠定的。因为他教导这种德性也作为一种保存者、维持者和守护者。\par

% structure: p0024 source=h2 target=section reason=relative-heading-rank; base=h2

\section{论据十二。—— \Person{国瑞}不许在自己方面获得任何德性。虔敬既是开始也是终结，因此是一切德性之母。}

诚然，由于我们迟钝而懒惰的本性，他尚未成功地使我们成为正义、明智、节制或勇敢的人，尽管他热切地在我们身上劳苦耕耘。因为实际上，我们并不拥有任何真正的德行，无论是人的还是神的，也从未接近过它们，反而仍然远离它们。这些都是非常伟大而崇高的德行，没有人可以平白获得，只有蒙天主赋予能力的人才能得到。我们的本性也绝非如此适合这些德行，我们也尚未声称自己配得上达到它们；因为由于我们的懈怠和软弱，我们没有做到那些追求最崇高、追求完美的人所当做的一切。因此我们还尚未成为正义的、节制的，或具有其他任何德行。但是这位可敬的人，这位德行的朋友和拥护者，早已为我们做了他力所能及的一切——使我们成为德行的热爱者，以最炽热的感情去热爱德行。而且，藉着他自己的德行，他在我们心中激发了对以下诸德的爱：正义之美（其金色面容实由他向我们展示）；明智，值得众人寻求；真正的智慧，最为甘饴；节制，这天上之德构成灵魂的健康状态，给拥有它的人带来平安；勇敢，这最可钦佩的恩惠；忍耐，这尤其属于我们的德行；而超越一切的是虔诚，人们正确地称它为德行之母。因为这是一切德行的开端和终结。始于这一德行，我们将发现其他所有德行最容易在我们身上增长：如果我们为自己热切渴求这恩宠（每一个只要不是彻底不虔或纯粹的享乐者的人都应当为自己获得它，以成为天主的朋友和祂真理的维护者），并且当我们勤勉地追求这一德行时，我们也留意其他德行，以便我们不致以不配和不洁的状态接近我们的天主，而是以一切德行和智慧作为我们最好的向导和最睿智的\OriginalTerm{司铎}{priests}。而我认为一切的终结无他，乃是：以纯洁的心灵相似天主，使你能接近祂并居住于祂内。\par

% structure: p0026 source=h2 target=section reason=relative-heading-rank; base=h2

\section{论题13。——\Person{奥力振}在他的神学和形而上学教导中所使用的方法。他劝勉研究所有作者，唯独无神论者除外。言谈中令人惊叹的说服力。心灵同意时的敏捷性。}

除了他其他一切忍耐而辛苦的努力外，我如何才能用言语述说他在教导我们神学和虔诚品格方面为我们所做的一切呢？我又如何能深入这位人物的真正心性，并表明他是以何等明智和审慎的准备，使我们熟悉一切关于天主性的论述，并谨慎地防范我们在最必要的事上（即对万有之因的认识）陷入任何危险呢？因为他认为我们应该以这样的方式研究哲学：极其勤奋地阅读所有由哲学家和古代诗人所写的著作，不拒绝任何东西，也不排斥任何东西（因为事实上我们还不具备批判性辨别的能力），唯独无神论者的作品除外，他们因自己的自负而背离人类的普遍理智，否认有天主或天主眷顾的存在。他要求我们远离这些作品，因为它们不值得阅读，并且因为我们内里那本该虔诚的灵魂可能会因听到违反敬拜天主的言辞而被玷污。因为即使是那些经常出入他们认为是虔诚殿堂的人，也小心不触碰任何亵渎之物。因此他认为，凡已承担起虔诚实践的人，根本不应考虑这些人的书。然而他认为，我们应该获取并熟悉所有其他著作，既不偏爱也不排斥任何一种，无论是哲学论述与否，无论是希腊文还是外邦文，而是倾听它们所有要传达的内容。他以极大的智慧和敏锐遵循这一原则，以免任何来自这个或那个学派的一句话被单独听取并被视为唯一真理，即使它可能并非真理，从而进入我们的心灵并欺骗我们，并且单独安顿在那里，使我们成为它的俘虏，以致我们再无力从中抽身，或洗净它，如同人洗掉一小块染了色的羊毛那样。因为人的言谈是强大而有力的，以其诡辩的言辞精妙，迅速进入耳中，塑造心灵，并用它所传达的内容印刻在我们身上；一旦它占据了我们的心，它就能赢得我们爱它如真理；即使它是虚假和欺骗的，它仍在我们的内里占据一席之地，如同某种魔法师般压倒我们，并保持其受骗的人作为它的拥护者。另一方面，人的心灵又是一种容易被言语欺骗的东西，非常容易同意；的确，在它正确辨别和探究问题之前，它很容易被说服，要么由于自身的迟钝和愚蠢，要么由于论说的精巧，而常常随意地、厌倦于精确考察，将自己交付给狡猾的推理和判断，这些判断本身是错误的，并且使接受它们的人陷入错误。不仅如此；如果另一种论说方式试图纠正它，它既不会接纳它，也不会让自己改变看法，因为它被先前占据它的任何想法牢牢抓住，仿佛有一个无情的暴君在统治它。\par

% structure: p0028 source=h2 target=section reason=relative-heading-rank; base=h2

\section{第十四论点。——哲学家的纷争源自何处。反对那些抓住一切遇见的东西就相信并紧抓不放的人。\Person{奥力振}习惯仔细阅读并向他的门徒解释外邦人的书籍。}

岂不是这样，互相矛盾和对立的学说被引入了，哲学家的所有争辩中，一派反对另一派的意见，一些人坚持某些立场，另一些人坚持其他立场，一个学派依附于一套教条，另一个学派依附于另一套？诚然，所有人都以哲学为志向，并声称自他们初次被唤醒以来就一直如此，还宣称他们如今在精通这些讨论之后，对它的渴望并不亚于开始时。不，他们反而声称，在仿佛浅尝了哲学并得以自由沉浸在讨论中之后，他们现在对哲学的爱甚至超过了当初，那时他们毫无经验，只是被某种冲动驱使去从事哲学。这就是他们所说的。从此以后，他们对任何持相反意见的人的话都不予理会。因此，古人中从未有人曾诱使任何现代人或逍遥学派信徒转向自己的思想方式，采纳自己的哲学方法；另一方面，现代人中也没有人将自己的观念强加于古学派的人。简而言之，没有人对其他人这样做过。因为要诱使一个人放弃自己的意见而接受他人的意见并非易事。尽管这些意见甚至可能会是，如果他在开始哲学化之前就相信它们，他最初本会完全乐意地钦佩和接受它们：因为当心灵尚未被占据时，他本会将注意力转向那套意见并予以赞同，并为了它们而反对他目前持有的意见。至少，我们高贵、最雄辩且精于批判的希腊人所展示的哲学正是这样的：因为无论这些人中的哪一个，在开始时被某种冲动所驱使而碰巧遇到的东西，他就独自宣称那是真理，并认为所有其他哲学家所坚持的一切都只是错觉和愚蠢，尽管他自己并不比其他人分别捍卫其特有教义更令人满意地通过论证来确立自己的立场。这个人的目的仅仅是不管通过强制还是劝说，都无需被迫放弃和改变自己的意见。而他（若可以说真话）除了某种对哲学方面这些教条的非理性冲动之外，别无所有，并且对于他认为真实的东西，除了（别惊异）不加辨别的偶然性之外，别无其他标准。由于每个人就这样依附于他最初遇到的立场，仿佛被锁链束缚，他不再能够注意其他意见。即使他碰巧对每个主题都有自己的见解并加以真理的证明，即使他有论证的帮助来显示其对手的教义是多么虚假，他也无能为力、漫不经心，仿佛期待某种偶然的巧合，他向最初占据他的那些推理屈服。这类推理不仅在其他事情上，而且首先在最重要和最根本的事情——在认识天主和虔敬上——误导那些接受它们的人。然而，人们如此被它们束缚，以致没有人能轻易释放他们。因为他们就像陷入一片广阔无垠的沼泽中的人，一旦跌入其中，既不能后退，也不能横跨过去以获救，而是被陷在泥泞中直到死亡。或者他们可以被比作进入一片深邃、浓密而壮丽的森林中的旅人，他或许以为能立即找到出路，重新走上开阔的平原，但被森林的广阔和茂密所挫败。他在其中转向各个方向，碰见各种连续的小径，走上许多道路，想着其中一条定能让他走出去。但这些路只会把他引向更深处，丝毫不能为他开辟出口，因为它们都只是森林内部的路。最后，这位旅人筋疲力尽，看到他尝试的所有路都只是森林，并因再也找不到地上的住所而绝望，决定留在那里，安家立业，在森林中开辟出他可使用的空地。或者，我们可以再取一个迷宫的比喻：它只有一个明显的入口，使人从外面想不到有什么诡计，于是通过那唯一的门进入。然后，走到最深处，观看那巧妙的景象，检查那构造得如此精巧、充满通道、以无尽的内外路径布置的设计，他决定出去，却发现自己无能为力，看到自己的出路完全被那看似聪明的内部构造所阻断。但毕竟，没有一个迷宫如此令人困惑和错综复杂，也没有一片森林如此浓密和曲折，也没有一片平原或沼泽能让那些一旦进入的人如此难以脱身，正如我们在某些这些哲学家的讨论中所遇到的那样。因此，为了避免我们陷入大多数人的不幸经历，他没有将我们引入任何单一的哲学学派；他也没有认为让我们只带着某一类哲学观点离开是合适的，而是将我们引入所有学派，并决定我们不应不知道任何希腊学说。他亲自与我们同行，在我们前面开路，并引领我们前行，如同在旅途中一般，每当遇到任何曲折、不健全和迷惑人的东西时。他像一位熟练的专家帮助我们，他对此类主题有长期的经验，对任何这类事情都不陌生，因而可以保持自身的高度安全，同时向他人伸出手，也确保他们的安全，如同将沉没者拉上来一样。他就是这样对待我们，从所有不同哲学家中选择并摆在我们面前一切有用和真实的东西，剔除一切虚假的东西。他在人类知识的其他分支中如此待我们，而尤其在一切关乎虔敬的事情上更是如此。\par

% structure: p0030 source=h2 target=section reason=relative-heading-rank; base=h2

\section{辩题十五：神圣事物之案例。唯天主及其先知当受听从于此。先知与听者受同一神感所动。\Person{奥力振}在解释\WorkTitle{圣经}上的卓越}

对于这些人间的导师，他确实劝诫我们不要依附于他们中的任何一个，即使他们被所有人见证为最有智慧的人，而只应专心于天主和先知。而他本人则成为我们解释先知的人，并阐明了其中一切晦暗或谜般的言辞。因为在神圣的话语中有许多这类事情；无论是天主乐意以这种方式与人沟通，以免圣言赤裸裸地进入一个不配的灵魂（像许多人的灵魂那样），还是因为每一个神圣的神谕就其本性而言是最清晰明了的，但对于我们这些背弃天主、因时间与年龄而失去听觉能力的人来说，似乎晦暗不明，我不得而知。但无论情况如何，如果确实有些言辞是谜一般的，他阐明了所有这一切，并将它们置于光明之中，因为他本人是一位熟练且最具辨别力的天主聆听者；或者如果它们中没有一个是真正晦暗的，那么它们对他也不是不可理解的，他是当今时代我所认识或从他人传闻中所知的所有人中，唯一如此深入研究天主清晰而光辉的神谕，以至于能够立即将它们的意义纳入自己的心灵，并传达给他人的。因为那位激励天主亲爱的先知、并启示他们所有预言及其神秘和属天话语的万民领袖，视此人为朋友，并立他为这些神谕的解释者；而他人仅得到暗示的事，他通过此人的工具得到了充分的教导；在一切可信赖的那一位或以君王方式命令、或简单宣告的事上，祂赐予此人探究、阐明和解释它们的恩赐：以便若有任何迟钝不信的人，或又渴望受教的人，能够从这人学习，并某种程度被强迫去理解、决定相信并跟随天主。此外，在我看来，这些事情他唯独且真实地通过参与圣神而发出：因为需要同样的能力给那些说预言的人和那些听先知的人；除非那说预言的同一圣神赐予他理解其话语的能力，否则没有人能正确地听先知。这一原则在\WorkTitle{圣经}本身中也得到了表达，\Quote{经上说：只有那关闭的才能打开，其他任何都不行；关闭的被打开，当受感的话语解释奥秘之时。}如今，这最大的恩赐此人从天主那里领受了，那最高贵的才干从天上赐予了他，使他成为天主神谕对人类的解释者，使他能理解天主的话语，如同天主对他自己说话一样，并能以世人可理解地聆听的方式向他们讲述。因此对我们来说，没有禁止谈论的话题；因为没有隐藏或无法接近的知识之事，但我们有能力学习每一种论述，无论是外邦的还是希腊的，精神的还是政治的，神圣的还是人类的；我们被允许以完全的自由走遍知识的全部领域，探究它，并以各种学说满足自己，享受智力的甘甜。无论摆在我们面前的是某个古老的真理体系，还是人们可能以其他方式命名的事物，我们在他身上拥有一个装置和力量，既令人赞叹又充满最美丽的观点。简言之，他真正是我们仿照天主乐园的一个乐园，在其中我们并非被安置去耕种脚下的泥土，或以身体营养使自己粗俗，而只是以一切喜乐和享受来增加心灵的收获——可以这么说，我们自己种植一些美好的生长物，或者由万有之源种植在我们里面。\par

% structure: p0032 source=h2 target=section reason=relative-heading-rank; base=h2

\section{辩题十六：\Person{额我略}以三重比喻哀叹他的离去；比拟为\Person{亚当}离开乐园、荡子离开父家，以及犹太人被掳往\Place{巴比伦}}

实在，这里就是\OriginalTerm{安慰的乐园}{paradise of comfort}；这里有真正的喜乐与欢愉，正如我们在这一时期所享受的，如今这时期即将结束——本身确非短暂，但若这真是结局，当我们离去、离开此地时，却又显得太短。因为我不知道是什么附了我，或我犯了什么过犯，使我如今要离去——如今要被赶走。我不知道该说什么，除非我就像\OriginalTerm{第二个亚当}{second Adam}，在乐园之外开始说话。我的生命何等美好，若我仅是聆听老师的讲论，而自己沉默！但愿如今我能学会缄默无言，而不是呈现这种新景象：使老师成为听者！我与这样的长篇大论有何相干？我有什么义务发表这样的演说，而理当不离此地、紧紧依附？但这些事似乎就如从最初欺骗而生的过犯，这些原始过犯的惩罚仍在此等待着我。我岂不于己视为悖逆，胆敢这样越过天主的话，而我本应停留在其中、紧守它们？我退出，逃离这种福乐生活，正如\OriginalTerm{原祖}{primeval man}躲避天主的面容，并回归我被取出的尘土。因此，我在那里将终生吃尘土，并要耕耘泥土——那给我长出荆棘和蒺藜的泥土，也就是说痛苦和令人责备的忧虑——我既脱离了美好而高贵的挂虑。我以前所离弃的，如今我返回——即那尘土，我由此而来，以及我在世上的寻常关系，和我父的家——离开那美好的土地，从前我不知那\OriginalTerm{美好的父家}{the good fatherland}所在；也离开那些关系，我在后期开始认出它们是我灵魂的真正亲属，并离开那为我们真正\Person{圣父}的家，\Person{圣父}居住其中，被真正的儿子虔敬尊崇，他们也渴望居住其中。但我，缺乏一切虔诚与配得，正从这些行列中出去，转身向后，循原路折回。经上记载，有一个儿子从他父亲那里领受了与另一位继承人、他兄弟相称的那份财物，便凭自己的决定，往他父亲极远的外方去了；在那里放荡生活，他把祖产挥霍散尽，完全浪费了；最后，因穷困所迫，他雇身为猪倌；被饥饿逼到绝境，他渴望分享喂猪的食物，却不能触到。如此他为自己放荡的生活付了代价，不得不以父亲那王侯般的餐桌，换得他未曾预料的——猪和奴婢的食物。我们似乎也面临某种类似的命运，如今我们离去，而且尚未领受我们所应得的全部分。因为虽然我们尚未领受所有应得的，我们仍然离去，留下与你们同在、在你们身边的尊贵与亲爱，只换取低劣之物。一切忧郁之事将陆续临到我们——骚动与混乱代替平安，混乱无规的生活代替宁静和谐的生活，残酷的束缚，市场的奴役，诉讼，人群，代替这份自由；既无欢乐，也无任何闲暇留给我们追求更高贵的目标。我们也不必谈论灵感的言语，却要谈论人的作品——先知视之为纯粹的祸害——而对我们来说，确实是恶人的作品。真的，我们将有黑夜代替白昼，黑暗代替明光，忧愁代替节庆的聚会；代替父家的，是敌国将要接纳我们，在那里我没有自由唱我的圣歌，因为在我灵魂陌生的土地上怎能唱它？在那里旅居者不得接近天主。只能哀哭悲伤，当我念及这里的迥异光景，如果那甚至在我能力之内。我们读到，敌人曾攻击一座伟大而神圣的城市，那里敬拜天主，他们把居民，无论歌者还是先知，都掳到他们自己的国家，即\Place{巴比伦}。且记载说这些俘虏被拘留在那地时，甚至当征服者要求他们唱圣歌或在异邦弹奏时，他们拒绝，将竖琴挂在柳树上，在\Place{巴比伦}河畔哭泣。我确实自视为他们中的一员，因我被逐出这城，离开我神圣的父家，在那里昼夜宣讲神圣律法，听见赞美诗、歌曲和属灵话语；那里有永久阳光；白天在醒时的异象中我们得以亲近天主的奥秘，夜间在梦中我们仍然专注灵魂在白天所见所触及的；简言之，那里神圣事物的灵感持续弥漫一切。我说，我从这城被逐出，被掳到异乡，在那里我无权吹笛：因为像这些古人，我将不得不把乐器挂在柳树上，河流将是我的寄居地，我将在泥土中劳作，无心唱圣歌，即使我记得它们；是的，可能因不断忙于其他事务，我甚至忘记它们，像一个被夺去记忆的人。但愿我离去时，仅是违背己愿，如俘虏常有的情形；但我也是自愿离去，并非受他人强制；我自己的行为使我失去这城，而我可以选择留在其中。或许离开此地，我也将走上不安全的旅程，如同有时离开某座平安之城的人；而确实很有可能在旅途中落入强盗之手，被擒，被剥去衣裳，被重重击伤，被抛弃在半死不活之地。\par

% structure: p0034 source=h2 target=section reason=relative-heading-rank; base=h2

\section{论点第十七——额我略自我安慰}

但我何必发出这样的悲叹呢？万民的救主仍然活着，甚至那半死和遭劫之人，也是所有人的保护者和医师，那圣言，是万物不眠的守护者。我们也拥有真理的种子，是您使我们知道那是我们的财富，以及我们从您那里领受的一切——那些宝贵的教导之存款，我们以此前行；尽管我们确实像离家远行的人一样流泪，但我们仍随身带着这些种子。那么，主管我们的守护者可能会安全地引领我们度过一切遭遇；我们也可能再次回到您身边，带来这些种子所结的果实和一把把收获，虽然远非完美——事实怎能如此？但仍是从事世俗事务的生活使我们可能培育的，固然被一种要么完全不能结果、要么倾向于结坏果实的才能所损坏，但我相信，这种才能不会注定被我们进一步滥用，如果天主赐予我们恩宠。\par

% structure: p0036 source=h2 target=section reason=relative-heading-rank; base=h2

\section{论点第十八——结束语，并为演说致歉}

因此，现在让我结束这席话。我竟敢在一个最不适合大胆的场合大胆发言。然而，我认为这篇讲话诚心诚意地尽我们所能表达了感谢——尽管我们没有什么配得上主题的话语，但我们也不能完全沉默——并且也表达了我们的遗憾，正如那些与朋友分离、远走他方的人常做的那样。我的这篇演说是否包含了幼稚 \Emph{或} 近乎阿谀奉承的内容，或是一方面过于简单、另一方面过于雕琢而令人反感，我不知道。但我清楚地意识到，其中至少没有不真实的东西，一切都是真实而纯粹的，出于真诚的意见，以及纯洁而完整的判断。\par

% structure: p0038 source=h2 target=section reason=relative-heading-rank; base=h2

\section{论点第十九——向奥力振呼求，并就此告别与迫切祈祷}

但是，亲爱的灵魂啊，起来祈祷，现在就释放我们吧；如同您以神圣的教导在您共融时拯救了我们，也请以您的祈祷在分离中依然拯救我们。请将我们交托，并时常将我们置于您的祈祷之中。或者，不断地将我们交托给那位带领我们到这时代的天主，为过去所赐的一切而感恩，并恳求祂将来继续引领我们，始终与我们同在，用对祂诫命的理解充满我们的心智，以对祂的敬畏激励我们，并赐予我们从今以后最好的引导。因为当我们离开您后，我们将不再有同与您在一起时那样的服从祂的自由。因此，请祈祷，当我们失去您的同在后，能从祂那里得到一些鼓励，并愿祂派遣一位好的向导，一位天使作为我们的同行者。也恳求祂转变我们的行程，因为这将是最能有效安慰我们、使我们重新回到您面前的一件事。\par

\SourceRecord{Translated by S.D.F. Salmond.；Ante-Nicene Fathers；Vol. 6.；Edited by Alexander Roberts, James Donaldson, and A. Cleveland Coxe.；Buffalo, NY: Christian Literature Publishing Co.,；1886.；<http://www.newadvent.org/fathers/0604.htm>.}
